%───────────────────────────────────────────────────────────────────────────────
\chapter{Tablas}
\label{ch:tablas}
%───────────────────────────────────────────────────────────────────────────────

Una tabla bien diseñada comunica estructura antes de que el lector lea un
solo número. El template construye tablas sobre tres principios del paquete
\texttt{booktabs}~\cite{booktabs2020}: nunca líneas verticales, reglas
horizontales con grosor diferenciado, y espacio generoso entre filas. El
coloreado de \texttt{catppuccin-palette.tex} se aplica como capa semántica
encima de esa estructura, no como sustituto de ella.

%───────────────────────────────────────────────────────────────────────────────
\section{Anatomía de una tabla}
\label{sec:tabla-anatomia}
%───────────────────────────────────────────────────────────────────────────────

\begin{definicion}{Reglas de booktabs}
\label{def:booktabs}
\texttt{booktabs}~\cite{booktabs2020} define tres reglas horizontales con
grosor decreciente. \texttt{\textbackslash toprule} abre la tabla (0.08em),
\texttt{\textbackslash midrule} separa el encabezado del cuerpo (0.05em),
y \texttt{\textbackslash bottomrule} la cierra (0.08em). Las reglas
intermedias parciales usan \texttt{\textbackslash cmidrule\{i-j\}}.
Nunca se usan líneas verticales (\texttt{|} en el preámbulo de columnas).
\end{definicion}

La estructura canónica de cualquier tabla del template es:

\begin{lstlisting}[style=latex, frame=none, numbers=none]
\begin{acadtable}[H]{tab:etiqueta}{Caption descriptivo}
  \begin{tabularx}{\textwidth}{l Y Z}
    \toprule
    \headerrow
    \thead{Col A} & \thead{Col B} & \thead{Col C} \\
    \midrule
    dato          & dato          & dato          \\
    \bottomrule
  \end{tabularx}
\end{acadtable}
\end{lstlisting}

Tres comandos propios del template:
\texttt{\textbackslash headerrow} aplica \texttt{\textbackslash rowcolor\{table-head\}}
a la fila del encabezado.
\texttt{\textbackslash thead\{\}} pone el texto en bold con color
\texttt{table-head-text}.
\texttt{\textbackslash tablenote\{a\}\{texto\}} añade una nota al pie.

%───────────────────────────────────────────────────────────────────────────────
\section{Tipos de columna}
\label{sec:tipos-columna}
%───────────────────────────────────────────────────────────────────────────────

\texttt{config/tables.tex} define tipos de columna complementarios a los
estándar de \LaTeX\ y \texttt{tabularx}~\cite{tabularx1999}:

\begin{acadtable}[H]{tab:col-tipos}{Tipos de columna disponibles}
  \begin{tabularx}{\textwidth}{l l X}
    \toprule
    \headerrow
    \thead{Tipo} & \thead{Definición} & \thead{Uso} \\
    \midrule
    \texttt{l}, \texttt{c}, \texttt{r}
      & Estándar \LaTeX
      & Ancho fijo por contenido, sin wrapping \\
    \texttt{X}
      & \texttt{tabularx}
      & Proporcional, izquierda, wrapping \\
    \texttt{Y}
      & \texttt{tabularx}
      & Proporcional, centrada, wrapping \\
    \texttt{Z}
      & \texttt{tabularx}
      & Proporcional, derecha, wrapping \\
    \texttt{L\{ancho\}}
      & \texttt{array}
      & Ancho fijo, izquierda, wrapping \\
    \texttt{C\{ancho\}}
      & \texttt{array}
      & Ancho fijo, centrada, wrapping \\
    \texttt{R\{ancho\}}
      & \texttt{array}
      & Ancho fijo, derecha, wrapping \\
    \bottomrule
  \end{tabularx}
\end{acadtable}

\begin{observacion}{Columnas proporcionales X/Y/Z}
Las columnas \texttt{X}, \texttt{Y}, \texttt{Z} solo funcionan dentro de
\texttt{tabularx}. Distribuyen el espacio sobrante de \texttt{\textbackslash textwidth}
(o \texttt{\textbackslash linewidth} dentro de tcolorbox) entre todas las
columnas de ese tipo de forma proporcional. Una tabla con columnas \texttt{lYY}
tiene una columna de ancho natural a la izquierda y dos columnas de igual
ancho proporcional. Todas usan centrado vertical (\texttt{m\{\}}) para
alinear celdas con distinto número de líneas.
\end{observacion}

%───────────────────────────────────────────────────────────────────────────────
\section{Entornos de tabla}
\label{sec:entornos-tabla}
%───────────────────────────────────────────────────────────────────────────────

\subsection{acadtable — tabla flotante}
\label{subsec:acadtable}

\texttt{acadtable} es un wrapper de \texttt{table} con caption y label
integrados. Es el entorno por defecto para tablas en flujo normal de texto:

\begin{lstlisting}[style=latex, frame=none, numbers=none]
% Placement por defecto: [htbp]
\begin{acadtable}{tab:label}{Caption del flotante}
  \begin{tabularx}{\textwidth}{...} ... \end{tabularx}
\end{acadtable}

% Placement forzado: usa [H] cuando la tabla es parte del argumento
\begin{acadtable}[H]{tab:label}{Caption}
  ...
\end{acadtable}
\end{lstlisting}

\subsection{boxtab — tabla no-flotante dentro de tcolorbox}
\label{subsec:boxtab}

Dentro de cualquier entorno basado en \texttt{tcolorbox} (\texttt{definicion},
\texttt{ejemplo}, etc.), los flotantes no están permitidos. El entorno
\texttt{boxtab} resuelve esto:

\begin{lstlisting}[style=latex, frame=none, numbers=none]
\begin{ejemplo}{Tabla dentro de un entorno}
  \begin{boxtab}{tab:label}{Caption de la tabla}
    \begin{tabularx}{\linewidth}{l Y}   % linewidth, no textwidth
      ...
    \end{tabularx}
  \end{boxtab}
\end{ejemplo}
\end{lstlisting}

\subsection{fittable — reescalado automático}
\label{subsec:fittable}

Para tablas con muchas columnas que exceden el ancho disponible,
\texttt{fittable} reescala el contenido hasta que quepa:

\begin{lstlisting}[style=latex, frame=none, numbers=none]
\begin{fittable}
  \begin{tabular}{llllllll}
    % tabla muy ancha que desbordaria la pagina
  \end{tabular}
\end{fittable}
\end{lstlisting}

\begin{observacion}{fittable y legibilidad}
\texttt{fittable} es un recurso de último recurso, no el primer intento.
Antes de usarlo, verificar si la tabla puede restructurarse: eliminar
columnas redundantes, dividir en dos tablas, o usar \texttt{longtable}
para tablas altas. Un texto reescalado al 60\% de su tamaño original
suele ser ilegible en impresión.
\end{observacion}

%───────────────────────────────────────────────────────────────────────────────
\section{Coloreado semántico}
\label{sec:tabla-color}
%───────────────────────────────────────────────────────────────────────────────

Los aliases del dominio \texttt{table-*} en la paleta controlan todos los
colores de tablas:

\begin{acadtable}[H]{tab:table-aliases}{Aliases de color para tablas}
  \begin{tabularx}{\textwidth}{l l X}
    \toprule
    \headerrow
    \thead{Alias} & \thead{Color Ctp} & \thead{Aplicación} \\
    \midrule
    \texttt{table-head}      & Blue!12   & Fondo del encabezado (\texttt{\textbackslash headerrow}) \\
    \texttt{table-head-text} & Text      & Texto del encabezado (\texttt{\textbackslash thead}) \\
    \texttt{table-row-alt}   & Surface0!30 & Filas alternadas (\texttt{\textbackslash rowcolor}) \\
    \texttt{table-rule}      & Overlay0!50 & Color de todas las reglas booktabs \\
    \bottomrule
  \end{tabularx}
\end{acadtable}

Para filas alternadas, \texttt{\textbackslash rowcolor\{table-row-alt\}} se
aplica manualmente en las filas pares. Cuando la tabla tiene muchas filas,
el paquete \texttt{colortbl} permite automatizarlo con
\texttt{\textbackslash rowcolors}:

\begin{lstlisting}[style=latex, frame=none, numbers=none]
% Alternar colores desde la fila 2 (la primera es el encabezado)
\rowcolors{2}{table-row-alt}{white}
\begin{tabularx}{\textwidth}{l Y}
  \toprule
  \headerrow
  \thead{Col A} & \thead{Col B} \\
  \midrule
  dato 1 & valor 1 \\
  dato 2 & valor 2 \\
  dato 3 & valor 3 \\
  \bottomrule
\end{tabularx}
\rowcolors{1}{}{}   % resetear para no afectar tablas posteriores
\end{lstlisting}

%───────────────────────────────────────────────────────────────────────────────
\section{Datos numéricos con siunitx}
\label{sec:tabla-siunitx}
%───────────────────────────────────────────────────────────────────────────────

\begin{definicion}{Columna \texttt{S} de siunitx}
\label{def:columna-s}
El paquete \texttt{siunitx}~\cite{siunitx2023} provee el tipo de columna
\texttt{S} que alinea números por el punto decimal, redondea al número de
dígitos especificado, y formatea incertidumbres. Es el tipo correcto para
cualquier columna de datos numéricos con unidades o cifras significativas.
\end{definicion}

\begin{acadtable}[H]{tab:convergencia}{Convergencia del método de Runge-Kutta 4 para $\equ{y}' = \eqop{-}\eqk{\lambda}\equ{y}$}
  \begin{tabularx}{\textwidth}{Y S[table-format=1.2e-1] S[table-format=1.4e-2] S[table-format=1.2]}
    \toprule
    \headerrow
    \multicolumn{1}{Y}{\thead{$\eqk{h}$}} &
    \multicolumn{1}{c}{\thead{Error $\|\equ{e}\|_\infty$}} &
    \multicolumn{1}{c}{\thead{Razón $\equ{e}_n/\equ{e}_{n-1}$}} &
    \multicolumn{1}{c}{\thead{Orden emp.}} \\
    \midrule
    $1/4$   & 2.31e-3  & \multicolumn{1}{c}{---} & \multicolumn{1}{c}{---} \\
    $1/8$   & 1.48e-4  & 0.0641 & 3.96 \\
    $1/16$  & 9.33e-6  & 0.0630 & 3.99 \\
    $1/32$  & 5.84e-7  & 0.0626 & 4.00 \\
    $1/64$  & 3.65e-8  & 0.0625 & 4.00 \\
    \bottomrule
  \end{tabularx}
  \tablenote{*}{$\eqk{\lambda} = 1$, $\equ{y}(0) = 1$, integración en $\eqdom{[0,1]}$.
  El orden empírico confirma la convergencia $\mathcal{O}(\eqk{h}^4)$ teórica.}
\end{acadtable}

\begin{observacion}{Celdas especiales en columna S}
Cuando una celda de una columna \texttt{S} contiene texto en lugar de un
número (como las entradas \texttt{---} de la primera fila), debe encerrarse
entre llaves: \texttt{\{---\}}. De otro modo \texttt{siunitx} intenta
parsearla como número y produce un error.
\end{observacion}

%───────────────────────────────────────────────────────────────────────────────
\section{Tablas con estructura compleja}
\label{sec:tabla-compleja}
%───────────────────────────────────────────────────────────────────────────────

\subsection{Celdas multifila y multicolumna}
\label{subsec:multirow}

Para tablas donde las categorías abarcan varias filas, el paquete
\texttt{multirow}~\cite{multirow2021} provee \texttt{\textbackslash multirow}
y \LaTeX\ provee \texttt{\textbackslash multicolumn}:

\begin{acadtable}[H]{tab:hiperparametros}{Comparación de hiperparámetros por arquitectura}
  \begin{tabularx}{\textwidth}{l l Y S[table-format=2.2] S[table-format=1.4]}
    \toprule
    \headerrow
    \thead{Modelo} & \thead{Optimizador} & \thead{Regularización} &
    \multicolumn{1}{c}{\thead{Acc. (\%)}} &
    \multicolumn{1}{c}{\thead{Loss}} \\
    \midrule
    \multirow{3}{*}{\textbf{ResNet-18}}
      & SGD       & Weight decay $10^{-4}$ & 93.12 & 0.2341 \\
      & Adam      & Dropout 0.3            & 92.87 & 0.2418 \\
      & AdamW     & Weight decay $10^{-3}$ & 93.54 & 0.2289 \\
    \midrule
    \multirow{2}{*}{\textbf{MobileNetV3}}
      & Adam      & Dropout 0.2            & 90.23 & 0.3102 \\
      & AdamW     & Weight decay $10^{-3}$ & 91.05 & 0.2987 \\
    \bottomrule
  \end{tabularx}
\end{acadtable}

\subsection{Cmidrule para separadores parciales}
\label{subsec:cmidrule}

\texttt{\textbackslash cmidrule\{i-j\}} dibuja una regla horizontal entre
las columnas $i$ y $j$, útil para separar grupos dentro de una tabla sin
usar \texttt{\textbackslash midrule} completo:

\begin{acadtable}[H]{tab:metodos-numericos}{Métodos numéricos para EDOs: orden y estabilidad}
  \begin{tabularx}{\textwidth}{l l C{2cm} X}
    \toprule
    \headerrow
    \thead{Familia} & \thead{Método} & \thead{Orden} & \thead{Observaciones} \\
    \midrule
    \multirow{3}{*}{Explícitos}
      & Euler         & 1 & Primer orden; inestable para $\eqk{\lambda h} > 2$ \\
      & Heun          & 2 & Predictor-corrector; región de estabilidad mayor \\
      & RK4           & 4 & Estándar en precisión; $4$ evaluaciones por paso \\
    \cmidrule{2-4}
    \multirow{2}{*}{Implícitos}
      & Euler implícito & 1 & A-estable; requiere resolver sistema no lineal \\
      & Crank-Nicolson  & 2 & A-estable; exacto para osciladores armónicos \\
    \cmidrule{2-4}
    \multirow{2}{*}{Multipaso}
      & Adams-Bashforth & 4 & Explícito; reutiliza evaluaciones anteriores \\
      & Adams-Moulton   & 5 & Implícito; par con Adams-Bashforth en PECE \\
    \bottomrule
  \end{tabularx}
\end{acadtable}

%───────────────────────────────────────────────────────────────────────────────
\section{Tablas dentro de entornos tcolorbox}
\label{sec:tabla-en-caja}
%───────────────────────────────────────────────────────────────────────────────

Como se explicó en la \cref{def:no-floats}, dentro de entornos
\texttt{tcolorbox} los flotantes no son permitidos. El entorno \texttt{boxtab}
y el ancho \texttt{\textbackslash linewidth} resuelven esto:

\begin{definicion}{Propiedades de la distribución normal multivariante}
\label{def:normal-mv}
Sea $\equ{\mathbf{x}} \sim \mathcal{N}(\eqk{\boldsymbol{\mu}},
\eqten{\boldsymbol{\Sigma}})$ con
$\eqten{\boldsymbol{\Sigma}} \in \eqdom{\mathbb{R}^{d\times d}}$ SPD.

\begin{boxtab}{tab:normal-props}{Propiedades principales}
  \begin{tabularx}{\linewidth}{l X}
    \toprule
    \headerrow
    \thead{Propiedad} & \thead{Expresión} \\
    \midrule
    Media
      & $\eqop{\mathbb{E}}[\equ{\mathbf{x}}] = \eqk{\boldsymbol{\mu}}$ \\
    Covarianza
      & $\eqop{\mathrm{Cov}}[\equ{\mathbf{x}}] = \eqten{\boldsymbol{\Sigma}}$ \\
    Entropía
      & $H = \tfrac{1}{2}\ln\det(2\eqcon{\pi e}\,\eqten{\boldsymbol{\Sigma}})$ \\
    Marginalización
      & Toda marginal es normal; $\equ{x}_i \sim \mathcal{N}(\eqk{\mu_i}, \eqten{\Sigma_{ii}})$ \\
    Transformación lineal
      & $\eqten{A}\equ{\mathbf{x}} + \eqk{\mathbf{b}} \sim
        \mathcal{N}(\eqten{A}\eqk{\boldsymbol{\mu}} + \eqk{\mathbf{b}},\,
        \eqten{A}\eqten{\boldsymbol{\Sigma}}\eqten{A}^\top)$ \\
    \bottomrule
  \end{tabularx}
\end{boxtab}

La densidad de probabilidad es
$\eqfn{p}(\equ{\mathbf{x}}) = (2\eqcon{\pi})^{-d/2}
\det(\eqten{\boldsymbol{\Sigma}})^{-1/2}
\exp\!\left(-\tfrac{1}{2}(\equ{\mathbf{x}}-\eqk{\boldsymbol{\mu}})^\top
\eqten{\boldsymbol{\Sigma}}^{-1}(\equ{\mathbf{x}}-\eqk{\boldsymbol{\mu}})\right)$.
\end{definicion}

%───────────────────────────────────────────────────────────────────────────────
\section{Importar datos desde CSV con pgfplotstable}
\label{sec:pgfplotstable}
%───────────────────────────────────────────────────────────────────────────────

\begin{definicion}{pgfplotstable}
\label{def:pgfplotstable}
El paquete \texttt{pgfplotstable}~\cite{pgfplotstable2023}, incluido en la
distribución de \texttt{pgfplots}, permite cargar archivos \texttt{.csv}
directamente y renderizarlos como tablas \LaTeX\ con formato automático
de números, redondeo y unidades.
\end{definicion}

\begin{lstlisting}[style=latex, frame=none, numbers=none]
% Cargar el archivo una vez (en el preambulo o en el capitulo)
\pgfplotstableread[col sep=comma]{data/resultados.csv}\datatable

% Renderizar como tabla con formato automatico
\pgfplotstabletypeset[
  columns        = {h, error, ratio, order},
  columns/h/.style      = {column name={$h$}},
  columns/error/.style  = {sci, sci e, precision=2,
                            column name={Error $\|\mathbf{e}\|_\infty$}},
  columns/ratio/.style  = {fixed, precision=4,
                            column name={Razon}},
  columns/order/.style  = {fixed, precision=2,
                            column name={Orden emp.}},
  every head row/.style = {before row=\toprule, after row=\midrule},
  every last row/.style = {after row=\bottomrule},
  col sep=comma,
]{\datatable}
\end{lstlisting}

\begin{observacion}{Formato del CSV}
\texttt{pgfplotstable} espera que la primera fila del CSV contenga los
nombres de las columnas sin espacios. Los números deben usar punto decimal
(no coma), independientemente de la configuración de \texttt{siunitx}.
El separador puede ser coma (\texttt{col sep=comma}), tabulador
(\texttt{col sep=tab}) o espacio (\texttt{col sep=space}).
\end{observacion}

%───────────────────────────────────────────────────────────────────────────────
\section{Nota al pie de tabla}
\label{sec:tablenote}
%───────────────────────────────────────────────────────────────────────────────

El comando \texttt{\textbackslash tablenote\{marca\}\{texto\}} añade una
nota inmediatamente debajo de la tabla. No usa \texttt{\textbackslash footnote}
(que envía la nota al pie de página) sino texto inline en tamaño
\texttt{\textbackslash footnotesize}:

\begin{lstlisting}[style=latex, frame=none, numbers=none]
\begin{acadtable}[H]{tab:label}{Caption}
  \begin{tabularx}{\textwidth}{l Y}
    ...
    \bottomrule
  \end{tabularx}
  \tablenote{*}{Resultados promediados sobre 5 semillas aleatorias.}
  \tablenote{\dag}{Modelo pre-entrenado en ImageNet; fine-tuning 50 epochs.}
\end{acadtable}
\end{lstlisting}

Varias llamadas a \texttt{\textbackslash tablenote} se apilan verticalmente.
La marca puede ser cualquier texto: \texttt{*}, \texttt{\dag},
\texttt{a}, \texttt{1}, etc.

%───────────────────────────────────────────────────────────────────────────────
\printbibliography[heading=subbibliography]
